% Options for packages loaded elsewhere
\PassOptionsToPackage{unicode}{hyperref}
\PassOptionsToPackage{hyphens}{url}
\documentclass[
]{article}
\usepackage{xcolor}
\usepackage[margin=1in]{geometry}
\usepackage{amsmath,amssymb}
\setcounter{secnumdepth}{-\maxdimen} % remove section numbering
\usepackage{iftex}
\ifPDFTeX
  \usepackage[T1]{fontenc}
  \usepackage[utf8]{inputenc}
  \usepackage{textcomp} % provide euro and other symbols
\else % if luatex or xetex
  \usepackage{unicode-math} % this also loads fontspec
  \defaultfontfeatures{Scale=MatchLowercase}
  \defaultfontfeatures[\rmfamily]{Ligatures=TeX,Scale=1}
\fi
\usepackage{lmodern}
\ifPDFTeX\else
  % xetex/luatex font selection
\fi
% Use upquote if available, for straight quotes in verbatim environments
\IfFileExists{upquote.sty}{\usepackage{upquote}}{}
\IfFileExists{microtype.sty}{% use microtype if available
  \usepackage[]{microtype}
  \UseMicrotypeSet[protrusion]{basicmath} % disable protrusion for tt fonts
}{}
\makeatletter
\@ifundefined{KOMAClassName}{% if non-KOMA class
  \IfFileExists{parskip.sty}{%
    \usepackage{parskip}
  }{% else
    \setlength{\parindent}{0pt}
    \setlength{\parskip}{6pt plus 2pt minus 1pt}}
}{% if KOMA class
  \KOMAoptions{parskip=half}}
\makeatother
\usepackage{color}
\usepackage{fancyvrb}
\newcommand{\VerbBar}{|}
\newcommand{\VERB}{\Verb[commandchars=\\\{\}]}
\DefineVerbatimEnvironment{Highlighting}{Verbatim}{commandchars=\\\{\}}
% Add ',fontsize=\small' for more characters per line
\usepackage{framed}
\definecolor{shadecolor}{RGB}{248,248,248}
\newenvironment{Shaded}{\begin{snugshade}}{\end{snugshade}}
\newcommand{\AlertTok}[1]{\textcolor[rgb]{0.94,0.16,0.16}{#1}}
\newcommand{\AnnotationTok}[1]{\textcolor[rgb]{0.56,0.35,0.01}{\textbf{\textit{#1}}}}
\newcommand{\AttributeTok}[1]{\textcolor[rgb]{0.13,0.29,0.53}{#1}}
\newcommand{\BaseNTok}[1]{\textcolor[rgb]{0.00,0.00,0.81}{#1}}
\newcommand{\BuiltInTok}[1]{#1}
\newcommand{\CharTok}[1]{\textcolor[rgb]{0.31,0.60,0.02}{#1}}
\newcommand{\CommentTok}[1]{\textcolor[rgb]{0.56,0.35,0.01}{\textit{#1}}}
\newcommand{\CommentVarTok}[1]{\textcolor[rgb]{0.56,0.35,0.01}{\textbf{\textit{#1}}}}
\newcommand{\ConstantTok}[1]{\textcolor[rgb]{0.56,0.35,0.01}{#1}}
\newcommand{\ControlFlowTok}[1]{\textcolor[rgb]{0.13,0.29,0.53}{\textbf{#1}}}
\newcommand{\DataTypeTok}[1]{\textcolor[rgb]{0.13,0.29,0.53}{#1}}
\newcommand{\DecValTok}[1]{\textcolor[rgb]{0.00,0.00,0.81}{#1}}
\newcommand{\DocumentationTok}[1]{\textcolor[rgb]{0.56,0.35,0.01}{\textbf{\textit{#1}}}}
\newcommand{\ErrorTok}[1]{\textcolor[rgb]{0.64,0.00,0.00}{\textbf{#1}}}
\newcommand{\ExtensionTok}[1]{#1}
\newcommand{\FloatTok}[1]{\textcolor[rgb]{0.00,0.00,0.81}{#1}}
\newcommand{\FunctionTok}[1]{\textcolor[rgb]{0.13,0.29,0.53}{\textbf{#1}}}
\newcommand{\ImportTok}[1]{#1}
\newcommand{\InformationTok}[1]{\textcolor[rgb]{0.56,0.35,0.01}{\textbf{\textit{#1}}}}
\newcommand{\KeywordTok}[1]{\textcolor[rgb]{0.13,0.29,0.53}{\textbf{#1}}}
\newcommand{\NormalTok}[1]{#1}
\newcommand{\OperatorTok}[1]{\textcolor[rgb]{0.81,0.36,0.00}{\textbf{#1}}}
\newcommand{\OtherTok}[1]{\textcolor[rgb]{0.56,0.35,0.01}{#1}}
\newcommand{\PreprocessorTok}[1]{\textcolor[rgb]{0.56,0.35,0.01}{\textit{#1}}}
\newcommand{\RegionMarkerTok}[1]{#1}
\newcommand{\SpecialCharTok}[1]{\textcolor[rgb]{0.81,0.36,0.00}{\textbf{#1}}}
\newcommand{\SpecialStringTok}[1]{\textcolor[rgb]{0.31,0.60,0.02}{#1}}
\newcommand{\StringTok}[1]{\textcolor[rgb]{0.31,0.60,0.02}{#1}}
\newcommand{\VariableTok}[1]{\textcolor[rgb]{0.00,0.00,0.00}{#1}}
\newcommand{\VerbatimStringTok}[1]{\textcolor[rgb]{0.31,0.60,0.02}{#1}}
\newcommand{\WarningTok}[1]{\textcolor[rgb]{0.56,0.35,0.01}{\textbf{\textit{#1}}}}
\usepackage{graphicx}
\makeatletter
\newsavebox\pandoc@box
\newcommand*\pandocbounded[1]{% scales image to fit in text height/width
  \sbox\pandoc@box{#1}%
  \Gscale@div\@tempa{\textheight}{\dimexpr\ht\pandoc@box+\dp\pandoc@box\relax}%
  \Gscale@div\@tempb{\linewidth}{\wd\pandoc@box}%
  \ifdim\@tempb\p@<\@tempa\p@\let\@tempa\@tempb\fi% select the smaller of both
  \ifdim\@tempa\p@<\p@\scalebox{\@tempa}{\usebox\pandoc@box}%
  \else\usebox{\pandoc@box}%
  \fi%
}
% Set default figure placement to htbp
\def\fps@figure{htbp}
\makeatother
\setlength{\emergencystretch}{3em} % prevent overfull lines
\providecommand{\tightlist}{%
  \setlength{\itemsep}{0pt}\setlength{\parskip}{0pt}}
\usepackage{bookmark}
\IfFileExists{xurl.sty}{\usepackage{xurl}}{} % add URL line breaks if available
\urlstyle{same}
\hypersetup{
  pdftitle={Bayesian Linear Regression},
  hidelinks,
  pdfcreator={LaTeX via pandoc}}

\title{Bayesian Linear Regression}
\author{}
\date{\vspace{-2.5em}2026-04-17}

\begin{document}
\maketitle

\subsection{Introduction}\label{introduction}

Bayesian linear regression keeps the same linear predictor and Normal
residual structure as ordinary least squares but treats every parameter
as a random variable with a prior. Inference returns a joint posterior
over coefficients and the residual scale rather than a point estimate
plus standard errors. The practical gain is direct: credible intervals
carry the probability statement that students often (incorrectly)
attribute to confidence intervals, predictions inherit propagated
uncertainty automatically, and small-sample analyses no longer rely on
asymptotic approximations of \(t\) distributions.

\subsection{Prerequisites}\label{prerequisites}

A working understanding of multiple linear regression, comfort with the
idea of a likelihood, and a basic sense of priors. A Stan-capable
backend (\texttt{rstan} or \texttt{cmdstanr}) is needed because
\texttt{brms} compiles models on demand; \texttt{rstanarm} ships with
pre-compiled models if compile time is a concern.

\subsection{Theory}\label{theory}

The model is

\[y_i = X_i \beta + \varepsilon_i, \qquad \varepsilon_i \sim \mathcal N(0, \sigma^2),\]

with priors \(\beta_j \sim \mathcal N(0, \tau_j)\) and
\(\sigma \sim \mathrm{Exponential}(\lambda)\) or a half-\(t\). With
weakly informative priors and reasonable sample size, the posterior mean
is close to the OLS estimate while small-\(n\) inference is regularised
toward the prior. The default sampler is Hamiltonian Monte Carlo with
the No-U-Turn termination criterion (NUTS); convergence is assessed
through \(\hat R\), bulk and tail effective sample size, and
divergent-transition counts.

\subsection{Assumptions}\label{assumptions}

Linearity in the linear predictor, additive Normal residuals with
constant variance, independence of observations, and proper priors.
Convergence diagnostics are part of the model assumptions in practice
--- a chain that has not mixed cannot deliver a credible posterior
summary.

\subsection{R Implementation}\label{r-implementation}

\begin{Shaded}
\begin{Highlighting}[]
\FunctionTok{library}\NormalTok{(brms)}

\FunctionTok{set.seed}\NormalTok{(}\DecValTok{2026}\NormalTok{)}
\NormalTok{d }\OtherTok{\textless{}{-}} \FunctionTok{data.frame}\NormalTok{(}\AttributeTok{x =} \FunctionTok{rnorm}\NormalTok{(}\DecValTok{100}\NormalTok{), }\AttributeTok{y =} \DecValTok{2} \SpecialCharTok{+} \FloatTok{1.5} \SpecialCharTok{*} \FunctionTok{rnorm}\NormalTok{(}\DecValTok{100}\NormalTok{))}

\NormalTok{fit }\OtherTok{\textless{}{-}} \FunctionTok{brm}\NormalTok{(y }\SpecialCharTok{\textasciitilde{}}\NormalTok{ x, }\AttributeTok{data =}\NormalTok{ d,}
           \AttributeTok{prior =} \FunctionTok{prior}\NormalTok{(}\FunctionTok{normal}\NormalTok{(}\DecValTok{0}\NormalTok{, }\FloatTok{2.5}\NormalTok{), }\AttributeTok{class =} \StringTok{"b"}\NormalTok{) }\SpecialCharTok{+}
                   \FunctionTok{prior}\NormalTok{(}\FunctionTok{normal}\NormalTok{(}\DecValTok{0}\NormalTok{, }\DecValTok{5}\NormalTok{),   }\AttributeTok{class =} \StringTok{"Intercept"}\NormalTok{) }\SpecialCharTok{+}
                   \FunctionTok{prior}\NormalTok{(}\FunctionTok{exponential}\NormalTok{(}\DecValTok{1}\NormalTok{), }\AttributeTok{class =} \StringTok{"sigma"}\NormalTok{),}
           \AttributeTok{chains =} \DecValTok{4}\NormalTok{, }\AttributeTok{iter =} \DecValTok{2000}\NormalTok{, }\AttributeTok{refresh =} \DecValTok{0}\NormalTok{)}

\FunctionTok{summary}\NormalTok{(fit)}
\FunctionTok{plot}\NormalTok{(fit)}
\end{Highlighting}
\end{Shaded}

\subsection{Output \& Results}\label{output-results}

\texttt{summary(fit)} returns posterior mean, standard error of the
mean, posterior standard deviation, 95 \% credible interval, \(\hat R\),
and bulk and tail effective sample size for every parameter.
\texttt{plot(fit)} shows trace plots (one line per chain) alongside
marginal density plots; well-mixed chains overlap throughout warm-up and
produce indistinguishable density plots.

\subsection{Interpretation}\label{interpretation}

A reporting sentence: ``A Bayesian linear regression with weakly
informative priors estimated the slope at 1.53 (95 \% CrI 1.23 to 1.84)
and the residual scale at 1.02 (95 \% CrI 0.90 to 1.17);
\(\hat R = 1.00\) for all parameters, with
\(n_{\text{eff}} > 3{,}500\).'' Always report a convergence statistic
alongside the interval --- without it, a reader cannot tell whether the
posterior summary is trustworthy.

\subsection{Practical Tips}\label{practical-tips}

\begin{itemize}
\tightlist
\item
  Standardise predictors before fitting so a single
  \(\mathcal N(0, 2.5)\) prior on slopes is sensible across columns.
\item
  Run a prior predictive check (\texttt{sample\_prior\ =\ "only"}) and
  look at the implied range of \texttt{y}; if it spans implausible
  values, tighten the prior before fitting to the data.
\item
  For interactions or polynomials, use \texttt{bs()} or \texttt{poly()}
  inside the formula but check that the prior scale still makes sense
  after the transformation.
\item
  Use \texttt{bayesplot::pp\_check(fit)} to overlay replicated datasets
  on the observed outcome; systematic deviation in the tails signals the
  wrong family.
\item
  Compare nested models with
  \texttt{loo\_compare(loo(fit1),\ loo(fit2))} rather than \(F\)-tests;
  the difference in expected log predictive density and its standard
  error is a richer summary than a \(p\)-value.
\item
  For non-Gaussian outcomes, change \texttt{family\ =\ gaussian} to
  \texttt{bernoulli()}, \texttt{poisson()}, \texttt{student()}, or any
  other supported family --- the prior structure carries over unchanged.
\end{itemize}

\subsection{R Packages Used}\label{r-packages-used}

\texttt{brms} for formula-based Bayesian regression, \texttt{rstan} or
\texttt{cmdstanr} as the sampling backend, \texttt{bayesplot} for
diagnostic and posterior predictive plots, and \texttt{loo} for
cross-validation.

\subsection{For Reviewers}\label{for-reviewers}

\textbf{What to look for in a paper using this method.}

\begin{itemize}
\tightlist
\item
  \textbf{Common misapplications.}
\item
  \textbf{Diagnostics that should be reported but often aren't.}
\item
  \textbf{Red flags in tables and figures.}
\item
  \textbf{What to verify.}
\item
  \textbf{What an adequate Methods paragraph must contain.}
\end{itemize}

\subsection{See also --- labs in this
chapter}\label{see-also-labs-in-this-chapter}

\begin{itemize}
\tightlist
\item
  \href{labs/lab-bayesian-thinking-control-vs-decision-error.Rmd}{Bayesian
  thinking; α control vs decision error}
\item
  \href{labs/lab-brms-stan-loo-hierarchical-models.Rmd}{brms / Stan,
  LOO, hierarchical models}
\end{itemize}

\end{document}
